% =====================================================================
\section{Fase de Exploración}
\label{sec:exploracion}

El objetivo de esta fase fue establecer qué se iba a construir, para quién y
con qué tecnologías, y planificar las iteraciones de la fase de producción.

\subsection{Establecimiento de los interesados}

\begin{table}[H]
\centering
\caption{Grupos de interesados del proyecto}
\label{tab:interesados}
\begin{tabularx}{\textwidth}{L{3.4cm} Y}
\toprule
\textbf{Interesado} & \textbf{Rol e interés} \\
\midrule
Desarrolladora (autora) & Responsable del análisis, diseño, implementación,
pruebas y documentación del sistema. \\
Tutor & Supervisión técnica del producto y de la memoria; validación de que el
sistema responde a los objetivos del perfil. \\
Docente de la asignatura & Seguimiento metodológico y verificación del
cumplimiento del cronograma y del formato. \\
Usuarios finales & Personas que entrenan en gimnasio, en los niveles
principiante, intermedio y avanzado (sección~\ref{sec:aa-descripcion}). \\
Administrador del catálogo & Rol técnico que carga y actualiza el catálogo de
máquinas fuera de la aplicación. \\
\bottomrule
\end{tabularx}
\end{table}

\subsection{Definición del alcance}

El alcance del producto se fijó a partir de los objetivos específicos y de las
limitaciones descritas en la sección~\ref{sec:alcance}. Dos decisiones de
alcance se tomaron en esta fase y condicionaron el resto del proyecto:

\begin{enumerate}
    \item \textbf{Inteligencia artificial basada en conocimiento para las
    rutinas.} Generar rutinas con un modelo de aprendizaje automático habría
    requerido un conjunto de datos de rutinas etiquetadas por resultados, del
    que no se dispone. En cambio, las pautas de prescripción del ejercicio
    están consolidadas en la literatura (sección~\ref{sec:mt-entrenamiento}).
    Por ello se optó por un sistema basado en reglas, que además produce
    recomendaciones explicables.
    \item \textbf{Realidad aumentada basada en cámara.} Se implementa la
    superposición de la ficha de la máquina sobre la imagen en vivo de la
    cámara. El reconocimiento automático con un clasificador de imágenes se
    diseña como un punto de extensión, porque requiere un conjunto de
    fotografías de máquinas con permiso de uso que no forma parte del
    alcance.
\end{enumerate}

\subsection{Requerimientos funcionales}

La Tabla~\ref{tab:rf} presenta los requerimientos funcionales y su relación
con los objetivos específicos.

\begin{longtable}{L{1.1cm} L{10.0cm} C{1.9cm}}
\caption{Requerimientos funcionales del sistema}
\label{tab:rf}\\
\toprule
\textbf{ID} & \textbf{Descripción} & \textbf{Objetivo} \\
\midrule
\endfirsthead
\toprule
\textbf{ID} & \textbf{Descripción} & \textbf{Objetivo} \\
\midrule
\endhead
\bottomrule
\endfoot
RF1 & Registrar usuarios e iniciar sesión con correo y contraseña o con una
cuenta de Google; recuperar la contraseña por correo y cerrar sesión. & Base \\
RF2 & Capturar en el primer ingreso el objetivo de entrenamiento y el nivel de
experiencia del usuario (\emph{onboarding}). & OE3 \\
RF3 & Consultar el catálogo de máquinas con búsqueda por texto (sin
distinguir tildes) y filtros por grupo muscular y por tipo de equipo. & OE1 \\
RF4 & Mostrar la ficha de cada máquina: descripción, músculos principales y
secundarios, guía de uso paso a paso, errores comunes, prevención de lesiones
y galería de variantes. & OE1 \\
RF5 & Mostrar los ejercicios asociados a cada máquina con su animación de
ejecución, el músculo objetivo y el nivel de dificultad. & OE2 \\
RF6 & Escanear una máquina con la cámara y superponer su ficha sobre la imagen
en vivo; si no se identifica automáticamente, permitir que el usuario la
confirme en una lista. & OE1 \\
RF7 & Generar automáticamente una rutina semanal a partir del objetivo y el
nivel del usuario. & OE3 \\
RF8 & Crear, editar y eliminar rutinas personalizadas organizadas por días;
agregar ejercicios desde la ficha de una máquina; marcar una rutina como
activa. & OE4 \\
RF9 & Calcular el índice de masa corporal, clasificarlo según la OMS y guardar
el historial de mediciones. & OE5 \\
RF10 & Registrar el progreso físico (peso y estatura) con cálculo automático
del IMC. & OE6 \\
RF11 & Visualizar un reporte del progreso con gráfico de la evolución del
peso, variación acumulada y categoría de IMC actual. & OE7 \\
RF12 & Consultar y editar los datos del perfil (nombre, género, estatura,
peso, objetivo y nivel). & Base \\
\end{longtable}

\subsection{Requerimientos no funcionales}

Los requerimientos no funcionales se organizaron según las características de
calidad de la norma ISO/IEC 25010 (Tabla~\ref{tab:rnf}).

\begin{table}[H]
\centering
\caption{Requerimientos no funcionales}
\label{tab:rnf}
\begin{tabularx}{\textwidth}{L{1.2cm} L{3.1cm} Y}
\toprule
\textbf{ID} & \textbf{Característica} & \textbf{Descripción} \\
\midrule
RNF1 & Portabilidad & La aplicación se ejecuta en dispositivos Android con
la versión mínima soportada por React Native~0.86 (Android~7.0, API~24). \\
RNF2 & Usabilidad & Interfaz en español, navegación por pestañas y acceso a
cualquier módulo principal en no más de dos toques desde el inicio. \\
RNF3 & Seguridad & Cada usuario solo puede leer y escribir sus propios datos;
el catálogo es de solo lectura; las reglas del servidor validan tipos y
rangos. \\
RNF4 & Fiabilidad & Los cambios en rutinas, IMC y progreso se reflejan en
tiempo real; la aplicación tolera la pérdida temporal de conexión gracias a la
caché de Firestore. \\
RNF5 & Eficiencia & El catálogo de máquinas se descarga una sola vez por
sesión y las URL de imágenes se almacenan en caché. \\
RNF6 & Mantenibilidad & Código en TypeScript estricto, organizado en capas
MVC, sin errores de compilación ni de análisis estático. \\
\bottomrule
\end{tabularx}
\end{table}

\subsection{Casos de uso}

La Figura~\ref{fig:casos-uso} presenta el diagrama de casos de uso del
sistema. El actor principal es el usuario; Firebase participa como actor
secundario que provee autenticación, persistencia y archivos.

\begin{figure}[H]
\centering
\begin{tikzpicture}[font=\footnotesize, lbl/.style={font=\scriptsize, fill=white, inner sep=1pt}]
  \actor{usr}{(0,-0.5)}{Usuario}
  \actor{fb}{(13.6,-0.5)}{Firebase}
  \node[draw, rounded corners, minimum width=9.8cm, minimum height=12.6cm,
        label={[font=\small\bfseries]above:Sistema \nombreApp}] (sis) at (6.75,-0.5) {};
  \node[casouso] (u1) at (4.2, 5.0) {Registrarse / iniciar sesión};
  \node[casouso] (u2) at (4.2, 3.8) {Completar onboarding};
  \node[casouso] (u3) at (4.2, 2.6) {Consultar guía de máquinas};
  \node[casouso] (u4) at (4.2, 1.4) {Ver ficha y ejercicios};
  \node[casouso] (u5) at (4.2, 0.2) {Escanear máquina (RA)};
  \node[casouso] (u6) at (4.2,-1.0) {Generar rutina automática};
  \node[casouso] (u7) at (4.2,-2.2) {Crear rutina personalizada};
  \node[casouso] (u8) at (4.2,-3.4) {Calcular IMC};
  \node[casouso] (u9) at (4.2,-4.6) {Registrar progreso};
  \node[casouso] (u11) at (4.2,-5.8) {Gestionar perfil};
  \node[casouso] (u12) at (9.3, 1.4) {Agregar ejercicio a rutina};
  \node[casouso] (u13) at (9.3,-1.0) {Activar rutina};
  \node[casouso] (u10) at (9.3,-4.6) {Ver reporte de progreso};
  \foreach \u in {u1,u2,u3,u4,u5,u6,u7,u8,u9,u11} \draw (usr) -- (\u.west);
  \draw[flechad] (u12) -- node[lbl, above=2pt]{«extend»} (u4);
  \draw[flechad] (u5) -- node[lbl, right=2pt]{«include»} (u4);
  \draw[flechad] (u13) -- node[lbl, above=2pt]{«extend»} (u6);
  \draw[flechad] (u10) -- node[lbl, above=2pt]{«extend»} (u9);
  \foreach \u in {u12,u13,u10} \draw (\u.east) -- (fb);
\end{tikzpicture}
\caption{Diagrama de casos de uso del sistema}
\label{fig:casos-uso}
\end{figure}

\subsection{Selección de tecnologías}

\subsubsection{Framework de desarrollo}

Se compararon tres alternativas para construir la aplicación Android
(Tabla~\ref{tab:comparativa-framework}). Se eligió React Native con Expo por
permitir una sola base de código tipada con TypeScript, por contar con módulos
mantenidos para la cámara y las imágenes animadas, y por la experiencia previa
de la autora con React y TypeScript, que reducía el riesgo del proyecto.

\begin{table}[H]
\centering
\caption{Comparativa de tecnologías de desarrollo móvil}
\label{tab:comparativa-framework}
\begin{tabularx}{\textwidth}{L{3.2cm} Y Y Y}
\toprule
\textbf{Criterio} & \textbf{Kotlin nativo} & \textbf{Flutter} &
\textbf{React Native + Expo} \\
\midrule
Lenguaje & Kotlin & Dart & TypeScript \\
Plataformas & Solo Android & Android, iOS, web & Android, iOS, web \\
Acceso a cámara & Nativo & \emph{Plugin} & \archivo{expo-camera} \\
Integración con Firebase & SDK oficial & \emph{Plugin} oficial & SDK web
modular \\
Curva de aprendizaje para la autora & Alta & Media & Baja \\
\midrule
\textbf{Decisión} & & & \textbf{Seleccionada} \\
\bottomrule
\end{tabularx}
\end{table}

\subsubsection{Servicios del servidor}

Para el servidor se evaluaron Firebase, Supabase y un servidor propio
(Tabla~\ref{tab:comparativa-backend}). Se eligió Firebase porque ofrece
persistencia sin conexión y escuchas en tiempo real integradas en el SDK,
almacenamiento de archivos, inicio de sesión con Google y un nivel gratuito
suficiente, sin necesidad de administrar infraestructura.

\begin{table}[H]
\centering
\caption{Comparativa de alternativas para el servidor}
\label{tab:comparativa-backend}
\begin{tabularx}{\textwidth}{L{3.2cm} Y Y Y}
\toprule
\textbf{Criterio} & \textbf{Firebase} & \textbf{Supabase} &
\textbf{Servidor propio} \\
\midrule
Modelo de datos & Documental (NoSQL) & Relacional (PostgreSQL) & A elección \\
Tiempo real & Integrado & Integrado & A implementar \\
Caché sin conexión & Integrada en el SDK & No integrada & A implementar \\
Almacenamiento de archivos & Cloud Storage & Storage & A implementar \\
Infraestructura a mantener & Ninguna & Ninguna & Servidor y base de datos \\
\midrule
\textbf{Decisión} & \textbf{Seleccionada} & & \\
\bottomrule
\end{tabularx}
\end{table}

\subsubsection{Fuentes de datos de ejercicios}

El contenido de la guía de máquinas y ejercicios requiere textos, imágenes y
animaciones. Se evaluaron las fuentes disponibles considerando su licencia,
ya que el producto es propiedad intelectual de la universidad
(Tabla~\ref{tab:fuentes-datos}).

\begin{table}[H]
\centering
\caption{Fuentes de datos de ejercicios evaluadas}
\label{tab:fuentes-datos}
\begin{tabularx}{\textwidth}{L{2.9cm} Y L{2.6cm} Y}
\toprule
\textbf{Fuente} & \textbf{Contenido} & \textbf{Licencia} & \textbf{Uso en el
proyecto} \\
\midrule
wger \citep{wger} & API REST con ejercicios, categorías y equipos. & Libre
(AGPL / CC) & Ejercicios para la generación de rutinas. \\
free-exercise-db \citep{freeexercisedb} & 876 ejercicios con imágenes de
posición inicial y final. & Dominio público & Imágenes y referencias de
ejercicios. \\
WorkoutX \citep{workoutx} & Ejercicios con animaciones GIF. & Plan gratuito
con condiciones de uso & 95 animaciones de ejecución, descargadas por la API
oficial. \\
Contenido propio & Descripciones, guías de uso, errores y prevención de
lesiones. & Propia & Texto de las fichas, contrastado con guías de
entrenamiento. \\
\bottomrule
\end{tabularx}
\end{table}

\subsection{Planificación de las iteraciones}

La funcionalidad se distribuyó en seis iteraciones de producción
(Tabla~\ref{tab:plan-iteraciones}), ordenadas de modo que cada una se apoye en
las anteriores: primero la identidad del usuario, luego el catálogo, que
alimenta al escáner y a las rutinas, y finalmente los módulos de seguimiento.

\begin{table}[H]
\centering
\caption{Plan de iteraciones de la fase de producción}
\label{tab:plan-iteraciones}
\begin{tabularx}{\textwidth}{C{1.0cm} Y L{3.0cm} C{2.0cm}}
\toprule
\textbf{Iter.} & \textbf{Módulo} & \textbf{Requerim.} &
\textbf{Objetivos} \\
\midrule
1 & Autenticación, onboarding y perfil & RF1, RF2, RF12 & Base, OE3 \\
2 & Guía de máquinas y ejercicios & RF3, RF4, RF5 & OE1, OE2 \\
3 & Escáner de máquinas con realidad aumentada & RF6 & OE1 \\
4 & Rutinas: generación automática y planificación & RF7, RF8 & OE3, OE4 \\
5 & Nutrición: calculadora de IMC & RF9 & OE5 \\
6 & Progreso físico y reportes & RF10, RF11 & OE6, OE7 \\
\bottomrule
\end{tabularx}
\end{table}

\subsection{Cronograma}

La Figura~\ref{fig:cronograma} muestra el cronograma del proyecto en
semanas, dentro del período de cinco a ocho meses que establece el Artículo~5
del Reglamento de Titulación.

\begin{figure}[H]
\centering
\begin{ganttchart}[
    hgrid, vgrid={*1{draw=black!10}},
    x unit=0.52cm, y unit chart=0.55cm,
    title label font=\scriptsize\bfseries,
    bar label font=\scriptsize,
    group label font=\scriptsize\bfseries,
    bar/.append style={fill=green!35, draw=black!60},
    group/.append style={fill=black!60},
    bar height=0.5, group peaks height=0.25
  ]{1}{24}
  \gantttitle{Semanas}{24} \\
  \gantttitlelist{1,...,24}{1} \\
  \ganttbar{Exploración}{1}{2} \\
  \ganttbar{Inicialización}{3}{4} \\
  \ganttgroup{Producción}{5}{18} \\
  \ganttbar{It.1 Autenticación}{5}{6} \\
  \ganttbar{It.2 Guía de máquinas}{7}{9} \\
  \ganttbar{It.3 Escáner RA}{10}{11} \\
  \ganttbar{It.4 Rutinas}{12}{14} \\
  \ganttbar{It.5 Nutrición}{15}{16} \\
  \ganttbar{It.6 Progreso}{17}{18} \\
  \ganttbar{Estabilización}{19}{20} \\
  \ganttbar{Pruebas del sistema}{21}{22} \\
  \ganttbar{Memoria y defensa}{1}{24}
\end{ganttchart}
\caption{Cronograma del proyecto}
\label{fig:cronograma}
\end{figure}
